% studentloan_core.tex
% Shared body for the six student-loan letters. BURDEN-SHIFT VERSION.
% Emphasis is on the FURNISHER's irreconcilable reporting, not on the consumer.
% Built only on the consumer's own credit reports (no NSLDS is disclosed).
% \input this file INSIDE \begin{document}, after the recipient/header block and the
% one-paragraph opening. Speaks of the furnisher in the third person ("the Department").

\subsection*{The Reported Date of First Delinquency Is Self-Contradictory on Its Face}

The Department reports a Date of First Delinquency (DOFD) of \textbf{10/06/2024} (Default Resolution Group tradelines) and \textbf{11/05/2024} (Nelnet tradelines). That date cannot be reconciled with the reporting of these same accounts across the Consumer's own credit reports:

\begin{itemize}[noitemsep]
  \item The \textbf{Equifax report dated 09/24/2024} reports these accounts as current / paid as agreed, ``payment deferred,'' with the \textbf{Date of First Delinquency field blank}.
  \item The \textbf{TransUnion report} prints an estimated removal date of \textbf{08/2023} on these accounts --- a removal date that could only have been generated from a DOFD no later than 2016.
  \item The \textbf{Equifax report dated 03/01/2026} then reports a DOFD of \textbf{10/06--11/05/2024}.
\end{itemize}

\textbf{A single loan portfolio cannot simultaneously carry a blank date of first delinquency, a 2016-era date (implied by the 08/2023 removal date), and a 2024 date. These reportings are mutually exclusive and, on their face, establish that the date of first delinquency has been altered.} The burden is on the Department, and on the reinvestigating bureau, to substantiate the DOFD actually being reported --- not merely to confirm that it was ``verified.''

\subsection*{A Date of First Delinquency May Advance Only Upon a Substantiated Payment}

Under 15 U.S.C.\ \S~1681c(c)(1), the DOFD is fixed at the commencement of the delinquency that leads to the derogatory status and does not advance unless the account is brought current by an actual payment. If the Department contends that a payment reset or advanced the DOFD to 2024, \textbf{it must produce proof of that payment.} It has furnished none. To the contrary, the Department's own 09/24/2024 reporting shows these accounts with \emph{no date of first delinquency at all} --- confirming that, as of that date, no qualifying delinquency had been established, let alone one in 2024. Absent documented proof of a payment that advanced the date, the 2024 DOFD is unsubstantiated, and the seven-year reporting period --- which runs from the true DOFD and is a hard limit under 15 U.S.C.\ \S~1681c(a) --- renders these tradelines obsolete.

\subsection*{The Reported ``Date of Last Payment'' Is Unsupported and Irreconcilable}

The Department furnished a ``Date of Last Payment'' of \textbf{08/06/2023}. That entry cannot be reconciled with the Department's own contemporaneous reporting of these same accounts as ``over 120 days past due'' and never brought current, and --- on the 09/24/2024 report --- with no date of first delinquency. The Department must substantiate the 08/06/2023 entry with documented proof of an actual payment, or correct and delete it. A payment date that cannot be substantiated, used to support a later DOFD, is a willful accuracy violation under 15 U.S.C.\ \S~1681s-2(a)(1)(A).

\subsection*{Monthly Re-Publication}

Per the Equifax Data Furnisher Guidebook, a furnisher must report its full file (entire portfolio) each month. The Department therefore re-publishes the disputed DOFD every reporting cycle; each monthly furnishing is a separate, ongoing inaccuracy.

\subsection*{Structural Mismatch: Ten Collection Tradelines, Nine Loans}

Nine (9) originating Direct Loans appear on the Consumer's reports, yet \textbf{ten (10)} Default Resolution Group collection tradelines are being reported. One collection tradeline has no corresponding originating loan. A collection tradeline that does not map to an originating obligation cannot be verified and must be deleted.

\subsection*{Failure of the Required Pre-Reporting Notice}

Master Promissory Note Paragraph 20 provides: ``We will notify you at least 30 days in advance that we plan to report default information to a consumer reporting agency\ldots You will be given a chance to ask for a review of the debt before we report it.'' Independently, 15 U.S.C.\ \S~1681s-2(a)(7) requires written notice of negative information no later than 30 days after furnishing it, and the Debt Collection Act, 31 U.S.C.\ \S~3711(e), requires pre-report notice for federal debts.

The Department's own furnished data forecloses any claim that these obligations were met. On the 03/01/2026 report the Department furnished \textbf{Metro 2 Narrative Code 118 (``Customer Unable to Locate Consumer'')} on each tradeline:
\begin{itemize}[noitemsep]
  \item If Code 118 is accurate, the Department could not have provided the required notice, or the ¶20 review opportunity, before reporting.
  \item If Code 118 is inaccurate --- the Consumer's address, cellular number, and email were current and on file throughout the relevant period --- the Department furnished false information.
\end{itemize}
The Consumer does not concede inability to be located; the Department is held to its own certification. Under either reading, the adverse information cannot stand as furnished.
